\section{Introducción}
\label{sec:introduccion}

\subsection{Contexto del problema}

El proceso de pre-sustentación de trabajos de titulación en la carrera de Ingeniería de Software de la
UTEQ involucra al menos cinco roles (estudiante, tutor, coordinador, jurado, presidente de tribunal) y
seis hitos secuenciales obligatorios: registro de solicitud, carga de anteproyecto, asignación de
jurados, programación de cronograma, evaluación por rúbrica y emisión de acta firmada. Previo a este
proyecto, ese flujo se gestionaba con formularios físicos/digitales sueltos y hojas de cálculo, sin un
sistema que garantizara la integridad del anteproyecto entregado, la trazabilidad de quién asignó a qué
jurado y cuándo, o el cálculo verificable de la nota ponderada final.

\textbf{Nota de honestidad metodológica sobre la cuantificación del problema:} la guía de este informe
exige un ``contexto cuantificado del problema''. El equipo no realizó una encuesta formal a la
coordinación académica de la carrera para obtener cifras institucionales verificadas (número de
pre-sustentaciones por periodo, tiempo promedio de gestión manual, tasa de errores administrativos). Esa
cuantificación institucional real no existe en este informe --- se declara la ausencia explícitamente en
vez de inventar cifras plausibles pero no verificables, lo cual sería exactamente el tipo de dato
fabricado que las reglas transversales de esta guía penalizan. Lo que sí se cuantifica, porque es
verificable directamente en el repositorio, es el alcance técnico del artefacto construido: 15
requisitos funcionales (12 verificados con prueba automatizada real, Sección~\ref{sec:requisitos}), 31
controladores REST, y la evidencia empírica completa de la Sección~\ref{sec:evaluacion}.

\subsection{Preguntas de investigación}

\begin{description}
    \item[RQ1] ¿Puede automatizarse el flujo completo de pre-sustentación (solicitud $\rightarrow$
    anteproyecto $\rightarrow$ jurado $\rightarrow$ cronograma $\rightarrow$ evaluación $\rightarrow$
    acta) en un sistema web único, con trazabilidad verificable de cada transición de estado?
    \item[RQ2] ¿Cuál es el comportamiento real de rendimiento del sistema bajo carga concurrente
    moderada (50 usuarios), y qué efecto cuantificable tiene una capa de caché (Redis) sobre la latencia
    de un endpoint que consume una API externa?
    \item[RQ3] ¿Qué vulnerabilidades de seguridad reales (no hipotéticas) existen en la implementación,
    y en qué medida las contramedidas aplicadas (JWT de 7 claims, rate limiting, cabeceras HTTP,
    consultas parametrizadas) las cierran, verificado con herramientas automatizadas independientes
    (find-sec-bugs, OWASP ZAP)?
    \item[RQ4] ¿Qué tan honesta y verificable puede mantenerse la documentación de un proyecto académico
    de este tipo bajo presión de plazos, dado que el propio proceso de este equipo incluyó episodios
    reales de datos fabricados que tuvieron que detectarse y corregirse?
\end{description}

\subsection{Objetivos}

\textbf{Objetivo general:} diseñar, implementar y evaluar empíricamente un sistema web que automatice el
proceso de pre-sustentación de trabajos de titulación de la UTEQ, con evidencia verificable y no
fabricada en cada etapa del ciclo de vida del software.

\textbf{Objetivos específicos:}
\begin{enumerate}[nosep]
    \item Especificar los requisitos del sistema bajo ISO/IEC/IEEE 29148:2018, con historias de usuario
    en formato Connextra e INVEST, y casos de uso bajo la plantilla de Cockburn (responde RQ1).
    \item Implementar el sistema con una arquitectura de tres capas (Angular, Spring Boot, PostgreSQL +
    Redis) documentada con el modelo C4 y registros de decisión arquitectónica (ADR) (responde RQ1).
    \item Medir empíricamente el rendimiento bajo carga con k6 y cuantificar el efecto de la caché con
    estadística inferencial no paramétrica (responde RQ2).
    \item Auditar la seguridad del sistema combinando revisión manual contra OWASP Top~10, análisis
    estático (SpotBugs/find-sec-bugs) y escaneo dinámico (OWASP ZAP) (responde RQ3).
    \item Documentar de forma trazable y verificable cada hallazgo del proceso, incluyendo los errores
    propios del equipo, en vez de presentar solo resultados favorables (responde RQ4).
\end{enumerate}

\subsection{Contribuciones}

\begin{enumerate}[nosep]
    \item Un sistema funcional de código abierto (licencia MIT) que automatiza el ciclo completo de
    pre-sustentación, con 90 requisitos trazados (64 funcionales y 26 no funcionales, SRS v1.0.1; 36 de
    50 requisitos Must verificados con prueba automatizada real, 72\,\%).
    \item Un conjunto de evidencia empírica real y reproducible (scripts, datos crudos, notebooks
    ejecutados) para rendimiento, seguridad, cobertura y calidad web, documentado en
    \texttt{docs/mediciones/DATA-PROVENANCE.md} con la procedencia de cada cifra citada.
    \item Un caso documentado, dentro del propio proceso de desarrollo, de detección y corrección de
    datos académicos fabricados (encuesta de usabilidad, métricas de rendimiento, citas de evidencia de
    trazabilidad) --- una contribución metodológica sobre integridad de la evidencia en proyectos
    académicos de ingeniería de software, más que un hallazgo técnico convencional.
    \item Una plantilla reutilizable de documentación de ingeniería de requisitos (SRS versionado,
    historias de usuario Connextra+Gherkin, casos de uso Cockburn, checklist INCOSE completo) aplicable a
    proyectos similares de la carrera.
\end{enumerate}

\subsection{Organización del documento}

El resto de este informe se organiza así: la Sección~\ref{sec:marco-teorico} presenta el marco teórico;
la Sección~\ref{sec:trabajos-relacionados}, los trabajos relacionados y la búsqueda de literatura
realizada; la Sección~\ref{sec:requisitos} resume la ingeniería de requisitos (documentada en detalle en
\texttt{docs/requisitos/}); la Sección~\ref{sec:metodos} instancia Design Science Research y GQM; la
Sección~\ref{sec:arquitectura} describe el diseño y la arquitectura; la Sección~\ref{sec:implementacion},
la implementación; la Sección~\ref{sec:evaluacion} presenta la evaluación empírica completa; las
Secciones~\ref{sec:discusion}--\ref{sec:conclusiones} discuten los resultados, las amenazas a la validez,
el trabajo futuro y las conclusiones; y la Sección final presenta las declaraciones obligatorias y las
referencias.
\newpage
